\begin{lemma}
For every real \(\delta>0\), there is \(n_0\) such that every \(n\ge n_0\)
satisfies
\[
1<\delta\sqrt{n\log n}.
\]
\end{lemma}
\begin{proof}
The natural numbers tend to infinity as real numbers.  Therefore
\(\log n\to\infty\), and in particular \(\log n\ge 1\) for all sufficiently
large \(n\).  On that eventual range, since \(n\ge0\), we have
\[
n=n\cdot 1\le n\log n.
\]
Because \(n\to\infty\), this lower bound shows \(n\log n\to\infty\).  The
square-root function tends to infinity at infinity, hence
\[
\sqrt{n\log n}\to\infty.
\]
Multiplication by the fixed positive constant \(\delta\) preserves divergence
to infinity, so
\[
\delta\sqrt{n\log n}\to\infty.
\]
Applying the definition of convergence to \(+\infty\) with bound \(1\) gives a
threshold \(n_0\) such that for every \(n\ge n_0\),
\[
1<\delta\sqrt{n\log n}.
\]
\end{proof}
